\SourceLeaf{442}%
bezeichnete, so gilt
\[
\begin{aligned}
\chi_1&=\chi+\operatorname{Area}L'L'_1P'_1P'\\
\psi'-\varphi'+A&=LL'_1L'L+\chi_1-\chi;
\end{aligned}
\]
allein
\[
\begin{aligned}
\varphi'&=\chi+B\\
\psi'&=\chi_1+B_1,
\end{aligned}
\]
also
\[
B_1-B+A=LL'_1L'L.
\]
Die Winkel des Dreiecks $LL'L'_1$ sind offenbar
\[
A,\qquad 180^{0}-B,\qquad B_1,
\]
also ihre Summe
\[
180^{0}+LL'_1L'L.
\]

Die Form des Beweises wird einiger Abänderung und Erläuterung bedürfen, wenn der Punkt $(3)$ innerhalb des Dreiecks liegt. Allgemein schließen wir aber:
\begin{quote}
»Die Summe der drei Winkel eines Dreiecks, das auf einer beliebigen gekrümmten Fläche durch kürzeste Linien gebildet wird, ist gleich der Summe von $180^{0}$ und dem Flächeninhalt des Dreiecks auf der Hilfskugel, dessen Rand durch die Punkte $L$ gebildet wird, die den Punkten auf dem Rand des ersten Dreiecks entsprechen. Der genannte Flächeninhalt des Dreiecks ist als positiv oder negativ anzusehen, je nachdem, ob sein Rand im selben Sinn wie die Figur oder im entgegengesetzten Sinn durchlaufen wird«.
\end{quote}
\noindent\hspace*{1.2em}Hieraus folgt ferner leicht, dass die Summe aller Winkel eines Polygons mit $n$ Seiten, die auf der gekrümmten Fläche kürzeste Linien sind, [gleich] der Summe von $(n-2)180^{0}+$ Flächeninhalt des Polygons auf der Kugeloberfläche ist, usw.

\Needspace{5\baselineskip}
\begin{center}16.\end{center}
\noindent\hspace*{1.2em}Wenn sich eine gekrümmte Fläche vollständig auf die andere abwickeln lässt, behalten offenbar alle Linien auf der ersten Fläche bei der Übertragung auf die andere ihre Länge, ebenso wie die Winkel, die durch das
